% Section 2 (restructure FRONT, research/paper-review/restructure/PLAN.md §1.4; DECISIONS §D2, §A8): related work in prose. The
% paragraphs keep the labels of the former remarks. rem:related:t5 moved to app-related.tex (W6), as the opening paragraph of
% app:related:t5; one sentence here points to it. New paragraph: two-part codes for density estimation (Barron and Cover, "studied by"
% only; no theorem attributed, DECISIONS §A8). Literature content at abstract level only. Former status and source records, as comments:
%   rem:related:posteriors \status{known; unverified; remark}\src{steeper \S1.3, remark after Prop 2.14}
%   rem:related:mdl        \status{known; unverified}\src{steeper \S1.3}
%   rem:related:cold       \status{known; remark}\src{steeper Rem 4.6}
%   rem:related:new        \status{remark}\src{steeper \S3.6; split/notes-final \S2.1 "Method and novelty", \S2.2 "Novelty"}
\section{Related work}\label{sec:related}\label{sec:related:work}

\paragraph{Generalised posteriors and two-part codes.}\label{rem:related:posteriors}
The tempered posterior $\Pi^{1/\lambda}$ is a generalised (Gibbs, fractional or safe) posterior with learning rate $1/\lambda$ \citep{zhang2006epsilon,bissiri2016general,bhattacharya2019fractional,grunwald2012safe,grunwald2017inconsistency}; its maximiser, the $\lambda$-MAP, is a two-part code with complexity weight $\lambda$ (\cref{def:setup:objectives}). Cold posteriors \citep{wenzel2020good}, which raise prior and likelihood to the same power, keep the ratio of $\lambda=1$.\label{rem:related:cold} For i.i.d.\ density estimation with a countable list of densities, two-part (MDL) estimators are studied by \citet{barron1991minimum}; \cref{thm:recovery:main}(a),(b) is the supervised analogue, for sub-probability kernels, an arbitrary input law and every schedule $\lambda_n=o(n)$, and it identifies the limit as the shortest $\mu$-almost everywhere equivalent programs. Its proof is the standard argument from Ville's inequality and the strong law, and we claim no novelty for it; what we add is the rate for fixed $\lambda>1$, the contrast with $\lambda=1$ and the necessity results (\cref{thm:recovery:main}(c),(d), \cref{prop:recovery:lambda1,prop:recovery:necessity}).

\paragraph{$\lambda$-weighted two-part MDL in noisy classification.}\label{rem:related:mdl}
The closest prior work studies two-part codes in noisy binary classification, where a hypothesis is a classifier together with a code for the label noise, essentially a product hypothesis. \citet{grunwald2007suboptimal} show that Bayes and MDL classifiers can behave suboptimally under misspecification. \citet{zhu2025quantifying} quantify overfitting along the regularisation path of two-part-code MDL in supervised classification, and \citet{zhu2026overfitting} study overfitting and generalisation of (PAC-)Bayesian prediction in noisy binary classification. We do not claim as new the qualitative finding that a bounded $\lambda$ fits the noise and that $\lambda$ must grow with $n$. Their target is the data distribution; for the teacher's errors ours deliberately differs from it (\cref{rem:limits:schedule}). An earlier report \citep[Section~8]{claude2026whatfollows} formalised the proposal for a general imitation channel; \cref{app:related:t5} lists what we re-derive from it.

\paragraph{Algorithmic statistics.}
The structure function and sophistication split the description of a string into a model and a data-to-model code \citep{koppel1987complexity,gacs2001algorithmic,vereshchagin2004kolmogorov,vereshchagin2017algorithmic}; restricting the model family, for example to Hamming balls \citep{vereshchagin2010rate}, changes the structure function. Product models, a parameter with private noise at each input, are such a restriction; \cref{sec:kinks} uses their structure function, and \cref{sec:parity} compares the two splits for a noisy parity.

\paragraph{Novelty.}\label{rem:related:new}
Known are $\lambda$-weighted MDL in noisy classification, two-part MDL for density estimation, the structure function, sophistication and restricted model families, generalised and tempered posteriors, and our tools (dominance bounds, Ville's inequality, the strong law, Kakutani's dichotomy, list decoding, Hoeffding's inequality). The hull--penalty relation itself is elementary convex geometry, and the method of \cref{thm:recovery:main}(a),(b) is standard. To our knowledge the following are new: the contrast under a steeper prior between classes with a universal member and the private-randomness class; the exact-recovery rates for fixed $\lambda>1$, with the contrast at $\lambda=1$ and the necessity results; the kink vocabulary, the converse and the flattening by a universal member; the product kink of the noisy parity against set models; the teacher theorem over all hypotheses through list decoding, and its extension to flip rates below $0.1233$; and the reuse factor.
